\section{Remote access}
\label{sec:remote-access}

The server binds inside the container's network namespace, which nothing
outside can reach. Two SSH forwards carry it out. Every hop terminates on
loopback, so no inference port is ever published to the office network.

\begin{lstlisting}
# 1. in the container: bind the bridge interface
gpu shell /shared/$USER/start-ollama.sh

# 2. on the developer VM: forward to the container
ssh -f -N -L 127.0.0.1:11434:172.17.0.4:11434 camimo

# 3. on the laptop: forward to the developer VM
ssh -N -L 11435:127.0.0.1:11434 eugen@192.168.1.57

# 4. verify, from a second terminal
curl http://localhost:11435/api/tags
\end{lstlisting}

The client then uses \cmd{http://localhost:11435/v1} as an OpenAI-compatible
base URL with any non-empty API key, or \cmd{/api/...} for the native
interface.

\medskip
\textbf{Security.} Ollama performs no authentication whatsoever. Whoever reaches
the port can run inference on the card and pull or delete models. Binding
\cmd{0.0.0.0} inside the container exposes it to every account on the GPU host
--- acceptable within a team, but a deliberate decision rather than a default.
Binding \cmd{0.0.0.0} on the developer VM would publish an unauthenticated
endpoint to the whole LAN and should not be done. Access from outside the
office requires a VPN; the developer VM has no public address.
